\section{Where direct decision fine-tuning fails}
\label{sec:boundary}

The same recipe that improves standard bounded decisions does not solve hard compositional ones, and on
that tier a frozen backbone with a few exemplars is not beaten by anything we trained. This section
reports that boundary, the rendering confound attached to every comparison in it, and one diagnosed
data defect that our first ablation was too small to confirm and a completed $2{\times}2$ establishes
at $p{=}5\mathrm{e}{-}10$.

\subsection{The hard compositional tier}
\label{sec:boundary:hard}

Breaking JevBench's hard tier down by sub-family (\Cref{fig:hard_families}, \Cref{app:hard-families})
localizes where training helps and where it does not. Held-out rubric styles, families, and grammars
\emph{inside} our own generators' distributions move sharply under workflow training
(\Cref{sec:train:mixture}, \Cref{app:decisionmix}), but almost none of it transfers to \emph{level-7}
composition -- decisions requiring temporal arithmetic, unit conversion, expected-value reasoning, or
multi-attribute trade-offs. On the 880-item level-7 set the uniform-guess rate is .441 (598 items at
$\numcands{=}2$, 224 at 3, 58 at 4; the majority-class rate is .250), and the checkpoints score .495
(\texttt{ts1b}), .544 (\texttt{tm1b}) and .620 (\texttt{tl1b}). Only the 14B is clearly clear of the
guess rate, and all three sit 25--40 points below the same models' scores on the held-out families,
grammars and styles \emph{inside} the generator's distribution (.84--.90). The gap, not the level, is
the finding: the curriculum that moves every in-distribution family by 30-plus points moves this one
by 5. It matches the JevBench hard tier's own temporal/probability/trade-off families.

Two readings are available and the per-family data does not support the one we initially favored.
\emph{Data coverage} would say our generators produce boolean rule composition up to six levels deep,
never produce the operations level 7 and JevBench's long\_policy/temporal/probability/tradeoff
sub-families require, and that no model acquires composition from data that does not contain it. That
account predicts the frozen backbone should fail these families too. It does not: on the same 111 hard
items the frozen 14B with three exemplars leads \texttt{tl1b} on \texttt{tradeoff} (.500 vs.\ .167),
\texttt{long\_policy} (.421 vs.\ .158), \texttt{ambiguous} (.714 vs.\ .429), \texttt{probability}
(.500 vs.\ .300) and \texttt{temporal\_numeric} (.333 vs.\ .200), and trails only on
\texttt{multi\_hop} (.444 vs.\ .611); the sub-family counts are single digits and no individual cell
separates, but the direction is uniform across every serial family
(\Cref{fig:hard_families}, \Cref{app:hard-families}). The backbone therefore already carries some of
the behavior these families need, and our fine-tuning does not preserve it --- which is the same
conclusion the aggregate paired test reaches in \Cref{sec:boundary:frozen}, and it is about training,
not about the representation being incapable in principle.

Missing coverage is still the best explanation for why \emph{training} did not add the behavior; it is
not an explanation for why training removed it. We cannot separate either from the alternative that a
single constant-depth pass is the binding constraint, and the control that would decide it --- the
same frozen backbone given a scratchpad on the same 111 items --- is one we have not run.
\Cref{sec:discussion} treats it as an open question rather than a conclusion.

\subsection{A frozen backbone with exemplars is not beaten on the hard tier}
\label{sec:boundary:frozen}

\Cref{tab:tl1b} compares the previous-recipe 14B checkpoint, the bundle that fixed the long-state
rendering, dropped over-length rows, added DecisionMix v2, the typed heads, and calibration-shaped
checkpoint selection, the same bundle with frozen-teacher distillation removed, and the frozen backbone
with three exemplars and no training at all.

\begin{table}[t]
  \centering
  \caption{14B comparisons. \textbf{JevBench: public 231-id subset, unranked}; 95\% cluster-bootstrap
  intervals over the paraphrase \texttt{group} field ($n_{\mathrm{eff}}{=}36$ standard states,
  $n{=}111$ hard items). \texttt{tl1b\_nokd} is identical to \texttt{tl1b} in every flag except
  \texttt{--distill\_beta 0}. Paired per-item bootstraps for the three comparisons that matter are in
  the text; marginal intervals overlap heavily and should not be differenced by eye.}
  \label{tab:tl1b}
  \small
  \setlength{\tabcolsep}{3pt}
  \begin{tabular}{lcccc}
    \toprule
    & \texttt{ladder\_14b} & \texttt{tl1b} & \texttt{tl1b\_nokd} & Frozen, 3-shot \\
    \midrule
    JevBench std & .875 [.792, .944] & .931 [.861, .986] & .917 [.833, .986] & .819 [.708, .917] \\
    JevBench hard & .468 [.378, .559] & .450 [.360, .541] & .477 [.387, .568] & \textbf{.559} [.468, .649] \\
    JevBench Brier std (hard) & .173 (.85) & .175 (.66) & .127 (--) & .30 (.60) \\
    \texttt{long\_policy} (hard subfamily) & .053 & .158 & .211 & \textbf{.421} \\
    Training val NLL & -- & .438 & .410 & -- \\
    Held-out score NLL & 2.87 & \textbf{0.95} & -- & -- \\
    Level-7 composition (guess rate .441) & -- & .620 & -- & -- \\
    \bottomrule
  \end{tabular}
\end{table}

\textbf{The frozen backbone leads the hard tier, and this survives a paired test.} A paired cluster
bootstrap on the per-item difference over the same 111 hard items gives frozen-minus-trained
$+.108$ $[+.027, +.189]$, $p{=}.015$ against \texttt{tl1b}. Against \texttt{tl1b\_nokd} the same test
gives $+.081$ $[-.009, +.171]$, $p{=}.082$ -- consistent in direction but not separable at this sample
size. So the result is real for the checkpoint we released and directionally consistent for its
matched sibling, and it should be stated at that resolution rather than as a uniform ``fine-tuning
hurts''. On the standard tier the ordering reverses (frozen-minus-trained $-.111$ $[-.250, +.014]$,
$p{=}.10$).

What this supports is that \emph{our current fine-tuning recipe does not preserve or improve the
backbone's hard-tier behavior}. It does not by itself show that direct readout cannot solve hard tasks,
that one forward pass is insufficient, or that autoregressive reasoning is required; the rendering
confound of \Cref{sec:boundary:rendering} weakens any such causal reading further, since we did not
sweep rendering for the frozen hard-tier control.

\textbf{Frozen-teacher distillation: no detectable effect either way.} A KL term against the
\emph{same} backbone's own frozen three-shot letter distribution is the one component of the 14B bundle
with its own matched control. On that pair, \texttt{tl1b\_nokd} is nominally ahead on hard accuracy
(.477 vs.\ .450), \texttt{long\_policy} (.211 vs.\ .158), training validation NLL (.410 vs.\ .438) and
standard-tier Brier (.127 vs.\ .175), while \texttt{tl1b} is nominally ahead on standard accuracy (.931
vs.\ .917). But the paired bootstraps are $-.027$ $[-.090, +.027]$ ($p{=}.45$) on hard and $+.014$
$[-.042, +.069]$ ($p{=}.81$) on standard. \textbf{No effect of same-backbone frozen-teacher
distillation is detectable at this sample size, in either direction.} The direction is consistent
across several metrics and it is inside the noise; we report it as an uninformative control, not as a
negative finding, and it is not evidence about distillation in general.

\textbf{The \texttt{ladder\_14b} $\to$ \texttt{tl1b} standard-tier gain is likewise not clean}
(.875 $\to$ .931; paired $+.056$ $[+.000, +.125]$, $p{=}.13$). The probability-quality change in the
same bundle is much larger relative to its own noise (held-out score NLL 2.87 $\to$ 0.95, hard-tier
Brier .85 $\to$ .66) and is the part of that result we stand behind. One consequence worth recording
for practitioners: our own pre-registered release gate for the large model was a point threshold
(``hard $\ge .559$'') on a quantity with a $\pm9$-point interval. A gate of that form is not decidable
at this sample size and should be restated on a lower confidence bound, or on a metric with more items.

\subsection{Rendering is a first-class experimental factor}
\label{sec:boundary:rendering}

How the state and the candidate set are rendered into the prompt is not a presentation detail. We have
two independent demonstrations of its size, from opposite directions.

\textbf{(i) Rendering alone, frozen model, no training.} Scoring the identical public JevBench items
with the identical frozen Qwen3.5 checkpoints, changing only the rendering to a replication of another
published system's format, moves standard-tier accuracy by up to 12.5 points
(\Cref{tab:rendering}). No weights, no exemplar count, and no item set differ between the two columns,
and the effect survives a paired per-item bootstrap at two of the three sizes.

\begin{table}[t]
  \centering
  \caption{Rendering alone, on frozen weights. Identical Qwen3.5 base checkpoints, identical public
  JevBench items (\textbf{public 231-id subset, unranked}), letter-logit readout in both columns; only
  the state/option rendering differs. $\Delta$ is a paired cluster bootstrap over the per-item
  difference on the same items, 20{,}000 resamples.}
  \label{tab:rendering}
  \small
  \begin{tabular}{llccc}
    \toprule
    Tier & Backbone & our rendering & replicated rendering & paired $\Delta$ [95\% CI], $p$ \\
    \midrule
    \multirow{3}{*}{standard} & Qwen3.5-2B & .583 & .597 & $+.014$ $[-.194, +.208]$, $p{=}.95$ \\
     & Qwen3.5-4B & .764 & .847 & $+.083$ $[+.014, +.167]$, $p{=}.048$ \\
     & Qwen3.5-9B & .806 & \textbf{.931} & $\mathbf{+.125}$ $[+.056, +.208]$, $p{<}.001$ \\
    \midrule
    \multirow{3}{*}{hard} & Qwen3.5-2B & .387 & .441 & $+.054$ $[-.045, +.153]$, $p{=}.33$ \\
     & Qwen3.5-4B & .495 & .468 & $-.027$ $[-.117, +.063]$, $p{=}.64$ \\
     & Qwen3.5-9B & .541 & .595 & $+.054$ $[-.027, +.135]$, $p{=}.22$ \\
    \bottomrule
  \end{tabular}
\end{table}

The largest of these swings, 12.5 points on a frozen 9B checkpoint, is of the same order as the entire
published standard-tier spread between leaderboard entries, and larger than any difference between our
own model variants in \Cref{tab:headline}. It is concentrated on the standard tier; the hard tier does
not move separably at any size.

\textbf{(ii) Train/test render mismatch, same items, no truncation.} On a purpose-built held-out
long-state evaluation set ($n{=}605$, no exact-state overlap with either training corpus, scored with a
4{,}096-token window so that nothing is truncated at evaluation time; state lengths $p50 = 1{,}965$
tokens, max 2{,}843), a model trained on facts-last renders scores \textbf{.830} overall when the
evaluation set is rendered facts-last and \textbf{.640} when the same 605 items are rendered
facts-first. On the $\numcands{=}2$ \texttt{policy\_permit} subset ($n{=}330$, majority-class floor
.612) the two renderings give .842 and .615 -- that is, moving the evidence block from after the request
to before it, with the evidence fully visible in both cases, costs 23 points and lands the model exactly
on its majority-class floor.

The implication is methodological and it applies to us first: \textbf{a benchmark score is a joint
measurement of model capability and interface/rendering compatibility}, so rendering must be controlled
when comparing architectures. We do not use this to dismiss any other system's results -- we cannot,
because we have not run their rendering on their weights under a controlled protocol. What we can say is
that a comparison between two decision architectures that does not hold rendering fixed is measuring
something other than the architectures, and that our own frozen-vs-trained comparison
(\Cref{sec:boundary:frozen}) inherits exactly this weakness. It is also why the leaderboard context for
JevBench lives in \Cref{app:leaderboard} rather than in the body.

\subsection{A verified data defect, a benchmark too small to see it, and a $2{\times}2$ that resolves it}
\label{sec:boundary:truncation}

The long-state generator originally rendered supporting facts \emph{after} the request text. Two facts
about this are exactly measured and not in doubt. First, the tokenizer's \texttt{truncation\_side} is
\texttt{'right'}, verified empirically: the window keeps the start of the render and drops the end.
Second, at the released 1{,}024-token training window this removed the supporting facts from
\textbf{98.8\%} of long-policy rows (100\% at the 256-token window used earlier in the project). Every
model trained on that version was optimized to answer long policies from states that no longer contained
the evidence for the answer. The state-length histogram is \Cref{fig:truncation}.

The natural next claim -- that this defect \emph{caused} the long-context collapse, and that fixing it
caused the recovery -- is the one we cannot make. The 14B run that fixed it changed five things at once.
We therefore ran the matched single-variable ablation: two 1.7B arms identical in every flag except
whether \texttt{data\_wf\_long} renders its facts at 0\% or 97.6\% of the state text (same rows, same
labels, matched state length, $p50 = 8{,}702$ characters in both arms), at \texttt{--max\_state 1024}
with \texttt{--drop\_truncated} off in both arms so that truncation bites exactly as it did at release.
\Cref{tab:trunc} reports it, and it is a null result.

\begin{table}[t]
  \centering
  \caption{Matched facts-first vs.\ facts-last training ablation, 1.7B, single variable.
  \textbf{JevBench: public 231-id subset, unranked}, 95\% cluster-bootstrap intervals. The
  \texttt{long\_policy} difference is in the predicted direction but does not separate from zero; the
  hard aggregate moves the other way.}
  \label{tab:trunc}
  \small
  \setlength{\tabcolsep}{4pt}
  \begin{tabular}{lccp{6.1cm}}
    \toprule
    & Arm A (facts-first) & Arm B (facts-last) & Test \\
    \midrule
    \texttt{long\_policy} & .316 (6/19) & .105 (2/19) & Fisher exact $p{=}.232$; bootstrap $\Delta$ $[-.053, +.474]$ \\
    JevBench hard & .378 [.288, .468] & .396 [.306, .486] & contradicts the predicted direction \\
    JevBench std & .750 [.625, .861] & .736 [.597, .861] & indistinguishable \\
    \bottomrule
  \end{tabular}
\end{table}

On the benchmark subfamily, then, the ablation is \emph{consistent with} the truncation mechanism,
\emph{underpowered to confirm it} at 19 items, and \emph{contradicted} by the hard-tier aggregate. Our
first reading of that outcome was that the mechanism could not be demonstrated and that the
render-mismatch effect of \Cref{sec:boundary:rendering}(ii) was what we had actually measured. The
completed design shows that reading was an artifact of the metric.

\textbf{The full $2{\times}2$: the truncation fix is real, and JevBench was the wrong instrument.} The
single-render comparison above confounds the variable under test with render compatibility, because
each arm was scored only in the render its own training used. We therefore scored \emph{both} arms on
the same 605 held-out long states under \emph{both} renders, with no truncation at evaluation
(\Cref{tab:trunc2x2}).

\begin{table}[t]
  \centering
  \caption{The completed $2{\times}2$. Both 1.7B arms of \Cref{tab:trunc}, scored on the same $n{=}605$
  held-out long states (4{,}096-token window, nothing truncated) under both renders. The diagonal is
  each model in its own matched condition, which removes the render-mismatch confound; the off-diagonal
  is the mismatch effect. Majority-class floors: \texttt{policy\_permit} .612 ($n{=}330$, $\numcands{=}2$),
  \texttt{action\_select} .233 ($n{=}275$, $\numcands{=}4$).}
  \label{tab:trunc2x2}
  \small
  \setlength{\tabcolsep}{5pt}
  \begin{tabular}{llrrr}
    \toprule
    Trained on & Scored on & Overall & \texttt{policy\_permit} & \texttt{action\_select} \\
    \midrule
    facts-first & facts-first \emph{(matched)} & \textbf{.942} & \textbf{.933} & \textbf{.953} \\
    facts-first & facts-last & .797 & .788 & .807 \\
    facts-last & facts-first & .640 & .615 \emph{(at floor)} & .669 \\
    facts-last & facts-last \emph{(matched)} & .830 & .842 & .815 \\
    \bottomrule
  \end{tabular}
\end{table}

Comparing each arm in its own matched condition, facts-first training is ahead by \textbf{11.2 points}
(.942 vs.\ .830, 95\% CI $[+.077, +.147]$, $z{=}6.2$, $p{=}5\mathrm{e}{-}10$; \texttt{policy\_permit}
$+.091$ $[+.043, +.139]$, \texttt{action\_select} $+.138$ $[+.086, +.190]$). The facts-first arm is also
the more robust of the two: switching the render against it costs 14.5 points, against 19.0 for the
facts-last arm.

Three effects were stacked on top of each other, and the $2{\times}2$ separates them. \textbf{(1)} The
truncation fix works, unambiguously at $n{=}605$. \textbf{(2)} Render mismatch is separately real and
large; it is what put the facts-last arm exactly on its majority-class floor (.615) in the single-render
comparison, and why that comparison read as a refutation. \textbf{(3)} JevBench's \texttt{long\_policy}
subfamily could not resolve either effect at $n{=}19$. We record the sequence rather than only the
conclusion, because the intermediate state of this experiment --- a pre-registered primary metric
returning $p{=}.232$ on a mechanism that a better-powered measurement puts at $p{=}5\mathrm{e}{-}10$ ---
is the failure mode we think is most worth publishing here: \emph{the metric was the problem, not the
mechanism}, and a single-render ablation cannot tell those apart by construction.

The transferable lessons are then two. A fixed-length window interacts with \emph{where} the
load-bearing content sits and not only with how long the render is, so a truncation-length ablation
alone would not surface this; and silent truncation can produce confidently wrong behavior by
systematically removing the evidence while preserving the labels
\citep{sambasivan2021datacascades}. Full distributions and per-family cells are in
\Cref{app:truncation}.
